\documentclass[12pt]{article}

\usepackage[margin=1in]{geometry}
\usepackage[utf8]{inputenc}
\usepackage[T1]{fontenc}
\usepackage{amsmath, amssymb, amsthm}
\usepackage{mathrsfs}
\usepackage{hyperref}
\usepackage{setspace}

\onehalfspacing

% Theorem environments
\newtheorem{theorem}{Theorem}[section]
\newtheorem{lemma}[theorem]{Lemma}
\newtheorem{corollary}[theorem]{Corollary}
\newtheorem{proposition}[theorem]{Proposition}
\theoremstyle{definition}
\newtheorem{definition}[theorem]{Definition}
\newtheorem{example}[theorem]{Example}
\theoremstyle{remark}
\newtheorem{remark}[theorem]{Remark}

% Shortcuts
\newcommand{\R}{\mathbb{R}}
\newcommand{\C}{\mathbb{C}}
\newcommand{\Z}{\mathbb{Z}}
\newcommand{\CP}{\mathbb{CP}}
\newcommand{\MU}{MU_*}
\newcommand{\Ustar}{U^*(pt)}
\newcommand{\del}{\partial}
\newcommand{\tensor}{\otimes}

\title{The Algebra of Manifolds:\\
       Quillen's Theorem on Cobordism and Formal Groups}
\author{Deepak Jassal\\
  \small University of Northern British Columbia\\
  \small MATH 420: Structure of Groups and Rings}
\date{April 5, 2026}

\begin{document}

\maketitle

\begin{abstract}
We survey the deep connection between the topology of smooth manifolds and the pure algebra of formal group laws, culminating in Quillen's 1969--1971 theorem identifying the complex cobordism ring with Lazard's universal ring. Beginning with foundational definitions of manifolds, vector bundles, homology, and cohomology, we construct the cobordism ring of stably complex manifolds and explain how tensor products of complex line bundles give rise to a formal group law over that ring. We then state Lazard's theorem---which shows purely algebraically that the universal formal group law lives over a polynomial ring---and explain how Quillen's theorem completes the picture by showing this algebraic universality is realized geometrically. The exposition aims to be accessible to readers with a background in abstract algebra and point-set topology.
\end{abstract}

\tableofcontents
\newpage

\section{Introduction}

One of the most striking discoveries in twentieth-century mathematics is that the topology of smooth manifolds---the geometry of curved spaces---is intimately controlled by purely algebraic structures called \emph{formal group laws}. This connection, made explicit by Daniel Quillen in his landmark 1969 announcement \cite{Qui69} and 1971 paper \cite{Qui71}, unified two fields that had been developing in parallel: the geometric theory of cobordism and the algebraic theory of formal groups.

The central result, now known as \emph{Quillen's theorem}, states that the ring of complex cobordism classes of stably complex manifolds is isomorphic, as a graded ring, to the Lazard ring $L$---the universal coefficient ring of the universal formal group law. This is remarkable: the Lazard ring was constructed by purely algebraic means \cite{Laz55}, with no reference to topology, yet it reappears as the cobordism ring of smooth manifolds.

This paper develops the ideas leading to Quillen's theorem. Section~\ref{sec:manifolds} recalls the definitions of smooth and complex manifolds. Section~\ref{sec:bundles} introduces vector bundles and Euler classes. Section~\ref{sec:homology} covers homology and cohomology. Section~\ref{sec:cobordism} constructs the cobordism ring and recalls the Milnor--Novikov structure theorem. Section~\ref{sec:fgl} introduces formal group laws and Lazard's theorem. Section~\ref{sec:bridge} explains how Euler classes of line bundles produce the formal group law of cobordism. Section~\ref{sec:quillen} states and interprets Quillen's theorem, and Section~\ref{sec:conclusion} concludes.

\section{Manifolds and Complex Manifolds}\label{sec:manifolds}

\subsection{Smooth Manifolds}

\begin{definition}
An $n$-dimensional smooth manifold is a Hausdorff, second-countable topological space $M$ equipped with a maximal smooth atlas: a collection of homeomorphisms (charts) $\varphi_\alpha : U_\alpha \to V_\alpha \subseteq \R^n$ from open sets $U_\alpha \subseteq M$ to open sets in $\R^n$, such that all transition maps $\varphi_\beta \circ \varphi_\alpha^{-1}$ are smooth wherever defined.
\end{definition}

Intuitively, a manifold is a space that locally looks like $\R^n$; the atlas records these local identifications \cite{Lee13}. Classical examples include the circle $S^1$, the $2$-sphere $S^2$, the torus $T^2 = S^1 \times S^1$, and the $n$-dimensional complex projective space $\CP^n$, which will play a central role as a generator of the cobordism ring.

The study of smooth manifolds from a global-topological perspective was revolutionized by Thom in the 1950s, who showed in \cite{Tho54} that global properties of manifolds could be studied via homotopy theory through his Thom complex construction. This laid the entire groundwork for the cobordism program.

\subsection{Complex Manifolds and Stably Complex Manifolds}

\begin{definition}
A complex manifold of complex dimension $n$ is a manifold whose charts take values in $\C^n \cong \R^{2n}$, with holomorphic transition maps.
\end{definition}

The rich interplay between the complex and real structures is central to cobordism theory. However, the most important class for our purposes is slightly larger:

\begin{definition}
A stably complex manifold is a smooth real manifold $M$ whose stable normal bundle---its normal bundle in a sufficiently high-dimensional Euclidean space---admits a complex structure.
\end{definition}

Every complex manifold is stably complex, but not conversely. The cobordism theory of stably complex manifolds is called \textbf{complex cobordism}, and its representing spectrum is denoted $MU$ \cite{WC25}.

\section{Vector Bundles and Euler Classes}\label{sec:bundles}

\subsection{Vector Bundles}

\begin{definition}
A vector bundle of rank $n$ over $\mathbb{F} \in \{\R, \C\}$ is a triple $(E, X, \pi)$ where $E$ (total space) and $X$ (base) are topological spaces, $\pi : E \to X$ is a continuous surjection, each fiber $\pi^{-1}(x)$ is an $\mathbb{F}$-vector space isomorphic to $\mathbb{F}^n$, and the bundle satisfies \emph{local triviality}: every $x \in X$ has a neighborhood $U$ with $\pi^{-1}(U) \cong U \times \mathbb{F}^n$ via a fiber-preserving homeomorphism.
\end{definition}

A \textbf{line bundle} is a vector bundle of rank $1$. The M\"obius strip is the archetypal non-trivial real line bundle over $S^1$: locally a product, but globally twisted. This illustrates the fundamental distinction between local and global structure that is central to the theory.

Every smooth manifold $M$ has a canonical vector bundle, its \textbf{tangent bundle} $TM$, whose fiber over a point $x$ is the tangent space $T_x M$. Dually, the \textbf{normal bundle} of a submanifold measures directions perpendicular to it.

\subsection{Euler Classes}

\begin{definition}
The Euler class $e(E) \in H^n(X; \Z)$ of a rank-$n$ real (or complex) vector bundle $E \to X$ is the obstruction to the existence of a nowhere-zero continuous section.
\end{definition}

Concretely, if $e(E) = 0$ then a nowhere-zero section exists (one can ``comb the hair'' on the bundle); if $e(E) \neq 0$ then every section must vanish somewhere. For a complex line bundle $L$, the Euler class coincides with the first Chern class: $e(L) = c_1(L) \in H^2(X; \Z)$.

The Euler class is the key link between geometry and formal group laws: it is the behavior of $e$ under tensor products of line bundles that produces the formal group law of cobordism.

\section{Homology and Cohomology}\label{sec:homology}

\subsection{Homology}

For a topological space $X$, the $n$-th singular homology group $H_n(X; \Z)$ measures the $n$-dimensional holes in $X$. Formally, it is defined as
\[
  H_n(X;\Z) = \frac{\{\text{$n$-cycles}\}}{\{\text{$n$-boundaries}\}},
\]
where an $n$-cycle is a formal $\Z$-linear combination of continuous maps $\sigma : \Delta^n \to X$ (singular simplices) with zero boundary, and an $n$-boundary is a cycle that bounds a higher-dimensional chain. The groups $H_0, H_1, H_2, \ldots$ record, respectively, the number of connected components, the non-contractible loops, the enclosed voids, and so on.

\begin{example}
For the torus $T^2$: $H_0(T^2) \cong \Z$, $H_1(T^2) \cong \Z^2$, $H_2(T^2) \cong \Z$. The two generators of $H_1$ correspond to the two independent loops on the torus, each of which cannot be shrunk to a point.
\end{example}

\subsection{Cohomology and the Cup Product}

The $n$-th cohomology group $H^n(X; \Z)$ is the $\Z$-linear dual of $H_n(X;\Z)$ (by the Universal Coefficient Theorem). Its crucial advantage over homology is the existence of the \textbf{cup product}
\[
  \smile \; : H^p(X;\Z) \times H^q(X;\Z) \longrightarrow H^{p+q}(X;\Z),
\]
which makes the graded group $H^*(X;\Z) = \bigoplus_{n \geq 0} H^n(X;\Z)$ into a graded commutative ring. This ring structure is essential for the definition and study of characteristic classes.

The \textbf{Chern classes} $c_i(E) \in H^{2i}(X;\Z)$ of a complex vector bundle $E \to X$ are canonical cohomology classes that classify the twisting of the bundle. For a line bundle $L$, the first Chern class $c_1(L) \in H^2(X;\Z)$ equals the Euler class $e(L)$, and the Chern classes of tensor products satisfy a formula that, in cobordism theory, becomes a formal group law.

\section{Cobordism}\label{sec:cobordism}

\subsection{The Cobordism Relation}

\begin{definition}
Two closed (compact, without boundary) smooth manifolds $M$ and $N$ of the same dimension are cobordant if there exists a compact manifold $W$ with boundary $\partial W = M \sqcup N$.
\end{definition}

Cobordism is an equivalence relation on manifolds: reflexivity holds because $M \times [0,1]$ is a cobordism from $M$ to itself; symmetry follows by reversing orientation; transitivity holds by gluing cobordisms along a shared boundary component.

\begin{example}
The following are cobordisms:
\begin{itemize}
  \item A circle $S^1$ is cobordant to the empty manifold, since $S^1 = \partial D^2$ (a circle bounds a disk).
  \item Two disjoint circles are cobordant to a single circle, via the ``pair of pants'' cobordism.
  \item The complex projective plane $\CP^2$ is not the boundary of any compact oriented $5$-manifold; its signature $\sigma(\CP^2) = 1 \neq 0$ prevents this.
\end{itemize}
\end{example}

Cobordism is a much coarser equivalence relation than diffeomorphism or homeomorphism; it is significantly easier to compute, and it is possible to classify manifolds up to cobordism in all dimensions \cite{WCob25}.

\subsection{The Complex Cobordism Ring}

Taking stably complex manifolds and forming equivalence classes under cobordism yields abelian groups. These assemble into the \textbf{complex cobordism ring}:
\[
  \MU = \bigoplus_{n \geq 0} MU_n,
\]
where $MU_n$ is the cobordism group of closed stably complex manifolds of real dimension $n$. Disjoint union gives the addition, Cartesian product gives the multiplication, the empty manifold is the zero element, and orientation reversal gives additive inverses (note $[M] + [\overline{M}] = [\partial(M \times [0,1])] = 0$).

The fundamental structure theorem of this ring was established independently by Milnor \cite{Mil60} and Novikov \cite{Nov60, Nov67}:

\begin{theorem}[Milnor--Novikov]\label{thm:milnor-novikov}
The complex cobordism ring is a polynomial ring over $\Z$ with one generator in each even real dimension:
\[
  \MU \cong \Z[x_2, x_4, x_6, \ldots], \qquad \deg x_{2n} = 2n.
\]
The generators may be taken to be the complex projective spaces: $x_{2n} = [\CP^n]$.
\end{theorem}

This theorem is proved using the Adams spectral sequence \cite{Adams58} and the structure theory of Hopf algebras; see \cite{Pan11} for a detailed account. The key point for our purposes is that $\MU$ is a free polynomial ring, generated by the complex projective spaces.

\section{Formal Group Laws}\label{sec:fgl}

\subsection{Definition and Examples}

\begin{definition}[{\cite{Laz55}}]
A commutative formal group law over a commutative ring $R$ is a formal power series $F(T_1, T_2) \in R[[T_1, T_2]]$ satisfying:
\begin{enumerate}
  \item \textbf{Identity:} $F(T, 0) = F(0, T) = T$,
  \item \textbf{Commutativity:} $F(T_1, T_2) = F(T_2, T_1)$,
  \item \textbf{Associativity:} $F(T_1, F(T_2, T_3)) = F(F(T_1, T_2), T_3)$.
\end{enumerate}
\end{definition}

The two simplest examples are the \textbf{additive} formal group law $F(T_1, T_2) = T_1 + T_2$ and the \textbf{multiplicative} formal group law $F(T_1, T_2) = T_1 + T_2 + T_1 T_2$. Both are commutative; the first comes from the additive group $\mathbb{G}_a$ and the second from the multiplicative group $\mathbb{G}_m$.

\begin{remark}
There is no need for an axiom requiring inverses: given $F$, one can always find a unique power series $G(T) \in R[[T]]$ with $F(T, G(T)) = 0$, generalizing the negation map. This follows from the implicit function theorem for formal power series \cite{Silverman09, WFG25}.
\end{remark}

A \textbf{homomorphism} of formal group laws is a power series $f(T)$ with $f(F(T_1, T_2)) = G(f(T_1), f(T_2))$; an \textbf{isomorphism} is a homomorphism with a compositional inverse; a \textbf{strict isomorphism} additionally satisfies $f(T) \equiv T \pmod{T^2}$.

\subsection{Lazard's Universal Ring}

The key theorem in formal group theory, due to Lazard \cite{Laz55}, is:

\begin{theorem}[Lazard]\label{thm:lazard}
\begin{enumerate}
  \item There exists a \textbf{universal} commutative formal group law $F_{\mathrm{univ}}$ over a ring $L$: for every commutative ring $R$ and every formal group law $F$ over $R$, there is a unique ring homomorphism $\phi : L \to R$ with $\phi_*(F_{\mathrm{univ}}) = F$.
  \item The ring $L$ is a polynomial algebra over $\Z$:
  \[
    L \cong \Z[t_1, t_2, t_3, \ldots], \qquad \deg t_i = 2i.
  \]
\end{enumerate}
\end{theorem}

Part (1) is nearly tautological: one constructs $L$ by writing $F(T_1, T_2) = \sum_{i,j} c_{i,j} T_1^i T_2^j$ with $c_{i,j}$ formal variables, and defines $L$ as the quotient of the polynomial ring $\Z[c_{i,j}]$ by the ideal generated by all relations forced by the group law axioms. The non-trivial content of the theorem is Part (2): that $L$ is actually a \emph{free} polynomial ring, with no unexpected relations among the $c_{i,j}$. This is proved using a careful analysis of symmetric 2-cocycles; see \cite{WL25} and \cite[Appendix~2]{Rav86} for detailed proofs.

\begin{remark}
Lazard's theorem is a statement of pure algebra. No topology is involved. The astonishing discovery of Quillen is that this algebraic universality is realized by the geometry of cobordism.
\end{remark}

\section{The Formal Group Law of Cobordism}\label{sec:bridge}

\subsection{Euler Classes and Tensor Products of Line Bundles}

The bridge between cobordism and formal group laws is provided by the behavior of Euler classes under tensor products of complex line bundles. Let $X$ be a stably complex manifold, and let $L_1, L_2 \to X$ be complex line bundles. In complex cobordism theory, the Euler class $e(L) \in U^2(X)$ is defined for every complex line bundle $L$. A computation using the projective bundle formula yields:

\begin{proposition}[{\cite{FGL12, Qui69}}]
There exists a formal power series $F_U(T_1, T_2) \in \Ustar[[T_1, T_2]]$ such that for any two complex line bundles $L_1, L_2$ over any space $X$,
\[
  e(L_1 \otimes L_2) = F_U\!\left(e(L_1),\, e(L_2)\right) \in U^2(X).
\]
Moreover, $F_U$ is a commutative formal group law over $\Ustar$.
\end{proposition}

The formal group law $F_U$ is called the \textbf{formal group law of complex cobordism}. The computation underlying the proposition uses the projective bundle formula \cite{Switzer75}. Its coefficients $a_{ij} \in U^{2(i+j-1)}(pt)$ are certain cobordism classes, expressible geometrically in terms of the Milnor hypersurfaces $H_{ij} \subset \CP^i \times \CP^j$ \cite{Milnor55, Rav13, Pan11}.

\subsection{The Classifying Map}

Since $\Ustar$ carries a formal group law $F_U$, Lazard's universality (Theorem~\ref{thm:lazard}) provides a unique ring homomorphism
\[
  \phi : L \longrightarrow \Ustar
\]
sending the universal generators $t_i \in L$ to the cobordism classes $a_{ij} \in \Ustar$, and sending $F_{\mathrm{univ}}$ to $F_U$. Lazard's theorem guarantees that $\phi$ exists and is unique, but does not imply it is an isomorphism. This is the content of Quillen's theorem.

\section{Quillen's Theorem}\label{sec:quillen}

\subsection{Statement}

In 1969 Quillen announced the following result \cite{Qui69}, and provided a complete proof in 1971 \cite{Qui71}:

\begin{theorem}[Quillen, 1969--1971]\label{thm:quillen}
The ring homomorphism $\phi : L \longrightarrow \Ustar$ is an isomorphism of graded rings. Consequently,
\[
  \Ustar \;\cong\; L \;\cong\; \Z[x_2, x_4, x_6, \ldots],
\]
where $x_{2n} \in U^{-2n}(pt)$ corresponds to $[\CP^n] \in MU_{2n}$ under the isomorphism $\Ustar \cong \MU$ given by Poincar\'e duality.
\end{theorem}

In other words, the formal group law of complex cobordism is the \emph{universal} formal group law; every other formal group law over any commutative ring is a pullback of $F_U$.

Quillen's proof \cite{Qui71} proceeds by elementary methods, using Steenrod operations to analyze the mod-$p$ reduction of the cobordism ring and verify that $\phi$ is an isomorphism by checking it is so rationally and at each prime separately. An earlier approach by Novikov \cite{Nov67} used the Adams spectral sequence. As Ravenel remarks in his survey \cite{Rav13}:

\begin{quote}
``In 1969 Quillen discovered a deep connection between complex cobordism and formal group laws which he announced in \cite{Qui69}. Algebraic topology has never been the same since.''
\end{quote}

\subsection{Summary of the Argument}

The logical flow of the proof can be summarized as follows:
\begin{enumerate}
  \item \textit{(Geometry)} Tensor products of complex line bundles produce, via Euler classes, a formal group law $F_U$ with coefficients in the cobordism ring $\Ustar$.
  \item \textit{(Algebra)} Lazard's theorem gives a canonical ring map $\phi : L \to \Ustar$ classifying $F_U$.
  \item \textit{(Computation)} One shows $\phi$ is an isomorphism: surjectivity follows from the fact that the Milnor hypersurfaces generate $\MU$ (by the Milnor--Novikov theorem, Theorem~\ref{thm:milnor-novikov}), and injectivity follows from comparing dimensions using Lazard's explicit polynomial description of $L$.
\end{enumerate}

The key insight is that both rings have the same simple structure (free polynomial on even-degree generators), and the natural map between them preserves this structure.

\section{Conclusion}\label{sec:conclusion}

Quillen's theorem stands as one of the great unifications in twentieth-century mathematics. Starting from the geometric notion of manifolds bounding one another, and the algebraic notion of a formal power series satisfying group-law axioms, the theorem reveals that these two worlds are identical: the most general formal group law is realized by the geometry of complex line bundles over stably complex manifolds.

The key ideas are: complex line bundles carry Euler classes that compose under tensor product via a formal group law; this formal group law lives over the cobordism ring $\Ustar$; Lazard's theorem provides a universal formal group law over a polynomial ring $L$; and Quillen showed the resulting map $L \to \Ustar$ is an isomorphism. The first three steps are topological; Lazard's theorem is pure algebra.

The implications rippled through algebraic topology for decades. The chromatic program, elliptic cohomology, topological modular forms, and the resolution of the Kervaire invariant problem all build on the foundation laid by Quillen's theorem. It remains a striking example of how abstract algebra illuminates the deepest structure of geometric spaces.

\bibliographystyle{amsplain}
\bibliography{references}

\end{document}